O \gls{ocr} é um subcampo da Visão Computacional e da Inteligência Artificial 
dedicado à conversão de imagens contendo texto (impresso, manuscrito ou gravado) 
em dados legíveis por máquina \cite{chen2021survey}. 
Originalmente voltado à digitalização de documentos, 
o escopo foi ampliado pelo \textit{Deep Learning} 
para aplicações industriais e de IoT \cite{jawas2018water}.

\subsubsubsection{Pipeline Tradicional}
Um sistema clássico de OCR opera em etapas interdependentes \cite{chen2021survey}:
\begin{enumerate}
    \item \textbf{Aquisição e pré-processamento}: melhoria de contraste, correção de inclinação (\textit{deskewing}), remoção de ruído e binarização; em dispositivos embarcados, etapas de segmentação semântica podem isolar a Região de Interesse (ROI) antes do reconhecimento.
    \item \textbf{Localização e segmentação de texto}: identificação das regiões de texto.
    \item \textbf{Reconhecimento de caracteres}: classificação dos segmentos ou da sequência completa.
    \item \textbf{Pós-processamento}: dicionários, correção ortográfica ou verificação sintática para mitigar erros residuais.
\end{enumerate}

Em OCR, a leitura de uma palavra ou sequência numérica se beneficia do contexto em ambas as direções: o caractere atual é mais bem interpretado quando se conhece tanto o que veio antes quanto o que veio depois. 
Por isso, a CRNN emprega tipicamente camadas \textbf{bidirecionais} de LSTM (BiLSTM), que processam a sequência da esquerda para a direita e da direita para a esquerda e concatenam as representações resultantes \cite{shi2017crnn, hannun2017ctc}.

\subsubsubsection{Evolução Tecnológica: dos Métodos Heurísticos ao Deep Learning}
A literatura descreve três gerações principais \cite{chen2021survey}:

\subparagraph{Primeira geração (Template Matching).}
Sobreposição pixel a pixel com gabaritos estáticos; 
alta vulnerabilidade a variações de escala, rotação e fontes.

\subparagraph{Segunda geração (Características manuais + classificadores clássicos).}
Uso de descritores como HOG e SIFT, 
combinados com SVM ou MLP; dependência de \textit{feature engineering}.

\subparagraph{Terceira geração (End-to-End Deep Learning).}
Elimina a segmentação explícita de caracteres; 
modelos \textit{end-to-end} aprendem simultaneamente representações 
visuais e dependência contextual sequencial \cite{shi2017crnn}.
